\section{Population-level signal: a provisional result}
\label{sec:genus}

Whether traits read from fitted models carry biological signal in aggregate
was tested by \emph{lot-blind nearest-neighbour classification}: each
specimen is assigned the genus of its nearest neighbour in trait space,
excluding specimens from the same collection lot so that nest-mates cannot
leak the answer, and the resulting accuracy is compared with a
\emph{permutation null} in which genus labels are shuffled. An earlier
full-corpus analysis reported accuracy well above this null. That corpus
could not be re-executed in the final computing environment, and the only
available re-check (22 specimens, 9 genera) was underpowered and is
directional at most. The result is therefore reported as \textbf{provisional}
and deliberately excluded from this report's headline claims; the protocol
details a reader would need to audit it (lot definition, number of
permutations, treatment of rare genera, scale normalisation, trait selection
before classification) must accompany the run artifacts before it is cited
as validated (Table~\ref{tab:evidence}).

What the design does show is why aggregate signal and individual reliability
are compatible (\S\ref{sec:scale-dependence}): averaging over many specimens
and vertices suppresses within-part localisation noise (\F{F3}), whereas the
rank of any single specimen against another does not enjoy that protection.
